% Folder 83, batch 09 (archivists' pages 161-180).
% Pass: Opus 5.5 (claude-opus-5-5), 2026-10-10 — first pass, unchecked against the pages by a human.
%
% All twenty pages are mathematics in his hand, a fast drafting hand whose
% formulas carry better than the connecting prose. Subject: quadratic
% coverings and quadratic divisors in V(L); affine lines with involution and
% their contracted product A * A; the torsor Q(A) of quadratic functions under
% L^{⊗2}; then (pp. 171-180) extensions 0 -> E° -> E -> F -> 0 with a
% lambda-section / lambda-splitting, their lambda-contracted product, and the
% question of a unit for it.

\documentclass[11pt,a4paper]{article}
\input{../preamble/grothendieck}

\folder{83}
\batch{9}
\pages{161}{180}
\foldertitle{[Icosaèdres, bi-icosaèdres et algèbres d'Azumaya de rang 4] : notes manuscrites (1982-1983, 1986, s.d.).}
\dating{1982-1986}
\watermark{Édition de démonstration}

\begin{document}

\page{161}
\note{la page s'ouvre sur une phrase complète ; l'argument (revêtements quadratiques, torseurs, produit $\ast$) vient des pages qui précèdent le lot.}
La \struck{catégorie} donnée d'un rev.\ quadratique $X$ de $S$, revient : la donnée d'un \struck{fibré inv.} $L$ inversible et d'un div.\ quadr.\ \struck{dans} $\check{\mathbb{V}}(L)$ (\uncertain{fini} sur $X$), ou encore : $L$ et la fonction de degré $\leq 2$ sur $\check{\mathbb{V}}(L)$, dont la partie homogène est $u \mapsto u^{\otimes 2}$.

\struck{(Dualement, cela revient} L'algèbre de $X$ \struck{dans} le dual \add{$H \simeq E\otimes L^{-1}$} de $E$, l'unité étant le dual de l'augmentation $E \to k$. La transformation par $\otimes L^{-1}$ de $q$ sur une forme quadratique sur $H$, (on sait la former \uncertain{déterminant}) égale à $1$ sur $\mathbf{1}$. Il y a sur $E$ une unique structure \struck{multiplication} d'alg.\ quadratique, pour laquelle $1$ est unité et $q_H$ est le déterminant. On récupère $\mathrm{Tr}$ \struck{comme} via la forme bil.\ associée à $q_H$ comme $x \mapsto \varphi_{q_H}(1,x)$, et dualement sur $L$, on récupère $t \in E$, \struck{comme} par la formule
\[
  \varphi_{q_E}(u, x) = u \otimes (x \wedge t) \in L^{\otimes 2},
  \qquad u \in L,\; x \in E,\; x\wedge t \in \det E \simeq L .
\]
\note{après $=$ il a d'abord écrit, puis biffé, un symbole devant $(x\wedge t)$, et devant $= L^{\otimes 2}$ un groupe de lettres biffé et illisible.}
\marginal{Attention aux conv.\ de signe !}

\page{162}
Ceci définit \add{(via $t$)} une involution sur $E$ compatible avec l'augmentation, i.e.\ une involution sur $A$ i.e.\ sur $\check{\mathbb{V}}(A)$, laquelle laisse le diviseur $X$ invariant et y induit l'involution canonique.

Attention, si on se donne d'avance une involution sur $L$, et cherche les div.\ quadr.\ sur $\check{\mathbb{V}}(L)$ qui \struck{sont stables} \add{la respectent, celles-ci étant stables par l'involution,} \struck{ce qui est par \ill{}}, mais cette condition n'est pas suffisante -- en car.\ 2, on peut prendre l'involution identique, alors qu'il y a des diviseurs quadratiques (\ill{}) pour lesquels l'involution n'est pas l'identité !) Mais \struck{car} 2 inv., cette difficulté disparaît -- cela se \uncertain{résume} : un $H$ div.\ sur $\check{\mathbb{V}}(L)$ \struck{$E$}, stable par $\lambda \mapsto -\lambda$, est bien l'image inverse par $L \xrightarrow{u \mapsto u^{\otimes 2}} L^{\otimes 2}$ d'un \struck{\ill{}} $\xi \in L^{2}$. (Il suffit $=$ que $2$ ne soit pas div.\ de zéro \struck{\ill{}} pour que cette dernière propriété soit vraie -- mais pour l'étendre aux $A$-affines quelconques il faut supposer que l'involution sur $A$ a un point fixe…)

\page{163}
donc l'application
\[
  \xi \longmapsto q_\xi, \qquad A \ast A \longrightarrow
  \left\{\text{formes quadratiques sur } E \text{ qui sont compatibles à } t\right\}
\]
est compatible avec l'application
\[
  A^{\otimes 2} \longrightarrow L^{\otimes 2}, \qquad \alpha \longmapsto -\alpha
\]
via les \struck{modules} \struck{d'opérateurs} \add{modules de translations}. Donc il \uncertain{sera vérifié} dès que notre convention, et \struck{pour} prendre $q_\xi$ tel que
\[
  \xi + q_\xi(x) = x \ast x .
\]
\note{après $=$ il a d'abord écrit puis biffé un $x$ et une étoile, avant $x \ast x$.}

Ainsi, on a un isom.\ de torseurs
\[
  \underbrace{\mathrm{Quadcomp}(E, \varepsilon, t)}_{\text{sur } (L,\sigma)} \simeq L \mathbin{\ast_\sigma} L
\]
moyennant quoi on les identifie les deux. Ceci posé, pour toute $q \in L \ast_\sigma L$, le diviseur quadratique associé n'est autre que l'image inverse de la section $q$ de $\check{\mathbb{V}}(L \ast_\sigma L)$ par $\mathbb{V}(L) \to \mathbb{V}(L \ast_\sigma L)$.
\note{« $\mathrm{Quadcomp}$ » : lecture d'une abréviation (« Quad.\ comp. ») pour les formes quadratiques compatibles.}

\struck{\ill{}} En termes de coordonnées, i.e.\ choisissant une origine dans $A$, de sorte que $A$ s'identifie à $L$, et $t$ est \struck{une} donné par $\tau \in L$,
\[
  \sigma_L(x) = \tau - x
\]
\note{la page s'arrête ici ; la suite du calcul en coordonnées ne se trouve pas en tête de la page suivante, qui reprend autrement.}

\page{164}
Cherchons les formes quadratiques (sur $E$) sur $E$ compatibles avec une involution \add{i.e.\ avec un $t \in E$ donné} déjà donnée, (Il est immédiat que ces formes forment un torseur sous l'\uncertain{un} des formes \add{(à val.\ dans $L^{\otimes 2}$)} qui se factorisent par $E/L \simeq k$, i.e.\ sous $L^{\otimes 2}$ lui-même. Or on a déjà un tel torseur, savoir $A \ast A$, et une application quadratique canonique
\[
  A \longrightarrow A \ast A
\]
\note{il a d'abord écrit, puis biffé, un second facteur après $A \ast$ à gauche de la flèche.}
(compatible avec chgt de base « dans » défini sur les schémas correspondants) savoir
\[
  x \longmapsto x \ast x .
\]
Ainsi, \struck{\ill{}} pour tout
\[
  \xi \in A \ast A,
\]
si on dénote par $q_\xi$ l'unique \struck{forme} \add{fonction} \add{polyn.\ de degré $\leq 2$} \struck{quadratique}
\[
  q_\xi : A \longrightarrow L^{\otimes 2}
\]
qui satisfait
\[
  \xi + q_\xi(x) = x \ast x \qquad \forall x \text{ dans } A
\]
on aura, pour $\alpha \in L^{\otimes 2}$, remplaçant $\xi$ par $\xi + \alpha$
\[
  (\xi + \alpha) + q_{\xi+\alpha}(x) = x \ast x
\]
donc
\[
  q_{\xi+\alpha}(x) = q_\xi(x) - \alpha
  \qquad \text{pour } x \in A, \text{ i.e. } \varepsilon(x) = 1
\]
\note{sous $-\alpha$ il écrit $\overset{\shortparallel}{} -\alpha\,\varepsilon(x)^2$ : $-\alpha$ y est égalé à $-\alpha\,\varepsilon(x)^{2}$.}

\page{165}
On peut construire un torseur
$Q = Q(L,\sigma)$ sous $L^{\otimes 2}$, et une fonction sous-quadratique
\[
  A \longrightarrow Q
\]
dont la partie homogène soit
\[
  u \longmapsto u^{\otimes 2} : L \longrightarrow L^{\otimes 2},
\]
et qui soit \emph{universelle} pour ces propriétés, \struck{P\ill{} ceci, qui les fonctions \ill{}} \add{de sorte que} les \uncertain{sous}-quadratiques à val.\ dans $L^{\otimes 2}$ comp.\ avec $\sigma$ soient simplement les compositions
\[
  A \longrightarrow Q \longrightarrow L^{\otimes 2}
\]
\note{sous la seconde flèche : « isom.\ des $L^{\otimes 2}$-tors.\ (correspondant au choix d'une origine) ».}
\note{devant $=$, dans $Q = Q(L,\sigma)$, une lettre biffée en indice de $Q$.}

Pour vérifier \ill{} que l'ens.\ des fonctions sous-quadr.\ (à valeurs dans $L^{\otimes 2}$) qui nous concernent, forment un torseur sous le module des formes quadr.\ qui se factorisent par $E/E^{\circ} \simeq k$, lequel module est isomorphe à $\mathrm{Quad}(k, L^{\otimes 2}) \simeq L^{\otimes 2}$. En termes

\page{166}
de fonctions, sous-quadratiques sur $A$, cela signifie simplement qu'elles forment un torseur sous $L^{\otimes 2}$, par addition de fonctions \emph{constantes}. Si $\Lambda$ est ce torseur, on a
\[
  \Lambda \times A \longrightarrow L^{\otimes 2}, \qquad (q, x) \longmapsto q(x)
\]
tel que
\[
  (q + \alpha, x) \longmapsto q(x) + \alpha \qquad \alpha \in L^{\otimes 2}
\]
et on prend
\[
  Q_\sigma = \mathrm{Hom}_{\text{torseurs inversés}}(\Lambda, L^{\otimes 2})
\]
d'où
\[
  \Lambda \simeq \mathrm{Hom}_{\text{torseurs sous } L^{\otimes 2}}(Q, L^{\otimes 2})
\]
d'où \struck{canoniquement} application can.
\[
  A \longrightarrow Q_\sigma
\]
\marginal{(c'est $x \mapsto x \otimes \sigma x \in \mathrm{Sym}^2(E) \simeq \mathrm{Quad}(\check{E}, k) \simeq \mathrm{Quad}(E, L^{\otimes 2})$, $\check{E} = E \otimes L^{-1}$)}
\note{la note entourée est dans la marge gauche, reliée à la flèche $A \to Q_\sigma$ ; au-dessus de « torseurs sous $L^{\otimes 2}$ », il a d'abord écrit « torseurs sur $L^{\otimes 2}$ ».}

\struck{Pour un choix} Le choix d'une $q$ \add{($\in \Lambda$)} \ill{} : celui d'une section de $L^{\otimes 2}$-torseurs i.e.\ d'une origine dans $Q_\sigma$ (car $Q_\sigma \simeq \Lambda$, mais cette bijection renverse l'op.\ de $L^{\otimes 2}$ en la symétrique $-$). Ceci dit, le diviseur quadratique \add{$X_q$} associé à une $q$, correspondant à une section $q$ de $Q_\sigma$,

\page{167}
n'est autre que l'image inverse de cette section.

\struck{On} On voit que $Q$ a bonne variance i.e.\ est celle des fonctions sous-quadratiques sur un torseur $A$ dans un autre $Q$, sous $L^{\otimes 2}$ cette fois, « compatibles » avec \struck{\ill{}} (\uncertain{un r.\ de} $L$ dans $L^{\otimes 2}$ \uncertain{et}) $t$ : ce sont \struck{les} les \add{fonctions} \struck{quadratiques}
\[
  q : A \longrightarrow Q
\]
telles que
\[
  q(x + \alpha) = q(x) \mathbin{\dot{-}} (t \wedge x) \otimes \alpha + \alpha \otimes \alpha ,
  \qquad \alpha \in L,\ t \wedge x \in L .
\]
\marginal{attention aux signes !}
\note{au-dessus du second $\otimes$, un signe repassé, peut-être un premier $\otimes$ biffé.}

\emph{Donc} telles fonctions \add{sous}-quadratiques \struck{quadratiques} \uncertain{affines} \add{\ill{}} forment un \uncertain{contexte}. Si on a une telle, cela fait $Q$ iso.\ à \ill{} \uncertain{puis} \ill{} une sol.\ des pbs universels. Et on vient de construire une telle \uncertain{à l'aide de} $L$. \struck{On peut} Ainsi, à toute « droite affine » $A$ \struck{a} \add{à} involution on a associé
\[
  Q(A, \sigma) = Q(A)
\]
et $A \longrightarrow Q(A)$ \uncertain{rev.\ de degré 2}.
\note{« $\mathbin{\dot{-}}$ » : il écrit un signe moins surmonté d'un point.}

\page{168}
Considérons une addition $(A_1, \sigma_1)$, $(A_2, \sigma_2)$ de droites affines \add{sous $L_1$}, \add{sous $L_2$} à involution, d'où
\[
  A = \text{\struck{$Q(A_1$}}\ A_1 \ast A_2 \qquad \text{tors.\ sous } L_1 \otimes L_2 .
\]
Je dis que
\[
  \underbrace{Q(A)}_{\text{sous } L^{\otimes 2}} \simeq
  \underbrace{Q(A_1)}_{\text{sous } L_1^{\otimes 2}} \ast
  \underbrace{Q(A_2)}_{\text{sous } L_2^{\otimes 2}}
  \qquad \text{(iso.\ canonique)}
\]
\marginal{??}
\note{dans la première ligne, $A = $ est suivi d'un « $Q(A_1$ » biffé, puis $A_1 \ast A_2$ ; dans $Q(A_1) \ast Q(A_2)$, l'étoile est écrite par-dessus un autre signe.}

Pour le voir \uncertain{nous} \ill{} : cette formule, se \uncertain{déduit} d'une définition sur une involution naturelle $Q(A)$ \ill{}.

En tout cas, \struck{en termes} on a une application comme \ill{} \uncertain{celle} :
\[
  Q(A_1) \times Q(A_2) \longrightarrow Q(A),
  \qquad q_1, q_2 \longmapsto q_1 \otimes q_2 \,|\, A_1 \ast A_2 \subset E_1 \otimes E_2
\]
\struck{donnée en termes de co\ill{}} qui s'explicite ainsi, si on choisit \struck{\ill{}} \add{pt base} $o_1 \in A_1$, $o_2 \in A_2$, de sorte que
\[
  A_i \simeq A_i^{\circ}, \qquad \sigma_i \text{ devient } x \longmapsto b_i - x
\]
les formes quadratiques sur $A_i$ \struck{compatibles} \add{compatibles} à $\sigma_i$ s'écrivent
\[
  q_{c_i} \text{ ou } q_{b_i, c_i} : x_i \longmapsto x_i^2 - b_i x_i + c_i ,
  \qquad c_i \in L^{\otimes 2} \text{ \uncertain{comme}}
\]
et on a alors
\[
  q_{c_1} \otimes q_{c_2} = q_c :
  x \longmapsto x^2 - (b_1 b_2) x + (b_2^2 c_1 + b_1^2 c_2 - 4 c_1 c_2) ,
  \qquad x \in L_1 \otimes L_2 .
\]
\note{la formule pour $q_{c_1}\otimes q_{c_2}$ est écrite en bas de page sans être contrôlée ; dans $-b_i x_i$, le $b_i$ est écrit par-dessus une autre lettre ; dans $b_2^2 c_1$, le $b_2$ est repassé.}

\page{169}
i.e.
\[
  q_{c_1} \otimes q_{c_2} = q_c \qquad c = b_2^2 c_1 + b_1^2 c_2 - 4 c_1 c_2
\]
Cela montre que pour $\alpha_1 \in L_1^{\otimes 2}$, on a
\[
  (q_1 + \alpha) \otimes q_2 = q_1 \otimes q_2 + \ddot{\alpha}\, \delta_2(q_2),
  \qquad \alpha \in L_1^{\otimes 2},\ \delta_2(q_2) \in L_2^{\otimes 2}
\]
(il suffit de faire $c \mapsto c + \alpha$ dans la formule pour $c$) et de même
\[
  q_1 \otimes (q_2 + \beta) = q_1 \otimes q_2 + \delta_1(q_1)\, \beta
\]
où
\[
  \delta_1(q_{c_1}) = b_1^2 - 4 c_1, \qquad \delta_2(q_{c_2}) = b_2^2 - 4 c_2 .
\]
\struck{Ainsi, il faut \uncertain{également}} De plus
\[
  \delta(q_1 \otimes q_2) = \delta(q_1)\, \delta(q_2),
  \qquad \delta(q_1) \in L_1^{\otimes 2},\ \delta(q_2) \in L_2^{\otimes 2},
\]
$\delta(q_1\otimes q_2)$ discriminant des formes sur modules \uncertain{sous-jacents} de $A_1 \ast A_2 = A$, dans $(L_1 \otimes L_2)^{\otimes 2}$.

On trouve \uncertain{diagramme} commutatif\supplied{]}

\begin{tikzcd}[column sep=large]
A_1 \times A_2 \arrow[r, "\ast"] \arrow[d, "\pi_1 \times \pi_2"'] & A = A_1 \ast A_2 \arrow[d, "\pi_A"] \\
Q_1 \times Q_2 \arrow[r, "\otimes"] \arrow[d, "\delta_{A_1} \times \delta_{A_2}"'] & Q = Q_{A, \sigma_A} \arrow[d, "\delta_A"] \\
L_1^{\otimes 2} \times L_2^{\otimes 2} \arrow[r, "{\text{can}, \sim}"] & (L_1 L_2)^{\otimes 2}
\end{tikzcd}

où, en termes des \struck{coor} choix des pts origines $o_1$ dans $A_1$, $o_2$ dans $A_2$, $o_1 \ast o_2$ dans $A = A_1 \ast A_2$, tout est explicité (identifiant $A_1$, $A_2$, $A_1 \ast A_2$ à $L_1$, $L_2$ et $L_1 \otimes L_2$ resp.) :
\note{il ouvre un crochet devant « On trouve » et le ferme après « commutatif » ; la flèche verticale de droite porte un $\pi_A$ repassé.}

\page{170}
\begin{align*}
  x_1 \ast x_2 &= (o_1 + x_1) \ast (o_2 + x_2)
    = o_1 \ast o_2 + \bigl(x_1 \otimes (\overbrace{\sigma_2(o_2) - o_2}^{b_2})\bigr) \\
  &\qquad + (\sigma_1(o_1) - o_1) \otimes x_2 - 2\, x_1 \otimes x_2 \\
  &= o + \underbrace{x_1 \otimes b_2 + b_1 \otimes x_2 - 2\, x_1 x_2}_{x_1 b_2 + b_1 x_2 - 2 x_1 x_2}
\end{align*}
$\pi_1 : A_1 \to Q$ transforme $x_1$ en la forme
\[
  \pi_1(x_1) = q_{x_1(b - x_1)} =
  \bigl(z_1 \mapsto z_1^2 - b_1 z_1 + x_1(b_1 - x_1) = (z_1 - x_1)(z_1 - (b_1 - x_1))\bigr)
\]
et de même pour $\pi_2$, donc $\pi_1(x_1) \otimes \pi_2(x_2)$ est la forme
\begin{align*}
  z \longmapsto{}& z^2 - b_1 b_2 z + b_2^2 x_1(b_1 - x_1) + b_1^2 x_2 (b_2 - x_2) \\
  &- 4 x_1 (b_1 - x_1) x_2 (b_2 - x_2)
\end{align*}
et
\begin{align*}
  \pi(x_1 \ast x_2) &= \bigl(z \mapsto z^2 - \underbrace{b_1 b_2}_{b} z
     + \underbrace{(x_1 \ast x_2)(b_1 b_2 - x_1 \ast x_2)}\bigr) \\
  &\qquad (x_1 b_2 + b_1 x_2 - 2 x_1 x_2)(b_1 b_2 - x_1 b_2 - b_1 x_2 + 2 x_1 x_2)
\end{align*}
\note{la seconde ligne développe le terme sous l'accolade ; dans la dernière parenthèse, le signe devant $b_1 x_2$ est repassé.}

La \uncertain{calcul} de l'égalité est un peu pénible, mais sa peut se faire (j'ai dû le simplifier précédemment par un autre biais, en choisissant $q_1$, $q_2$ et faisant la vérification en choisissant $o_i = x_i$ \struck{tel que $q_i(o_i) = 0$} i.e.\ il suffit de la faire pour $x_1 = 0$, $x_2 = 0$ ce qui est trivial.

On voudrait écrire \struck{$Q(A) = Q(A$}
\[
  Q(A_1 \ast A_2) \simeq Q(A_1) \ast Q(A_2)
\]
mais le second membre \ill{} \struck{\ill{}} \uncertain{distributive} \uncertain{principalement} un \ill{} n'y prête pas, \struck{\ill{}} \uncertain{les} $Q(A_i)$ ne sont pas \add{à} involution. Il faudrait reprendre en \ill{} \struck{\ill{}} la formalisation \add{dans} \uncertain{un cadre plus général} !
\note{les dernières lignes de la page, serrées en bas du feuillet, sont d'une lecture très incertaine.}

\page{171}
Considérons une extension
\[
  0 \longrightarrow E^{\circ} \longrightarrow E \xrightarrow{\ \varepsilon\ } F \longrightarrow 0,
  \qquad F \simeq E/E^{\circ}
\]
\uncertain{nous} allons déterminer les involutions de $E$, induisant $-\mathrm{id}_{E^{\circ}}$ sur $E^{\circ}$, et $+\mathrm{id}_F$ sur $F$. Une involution $\sigma$ de $E$ \uncertain{y} \ill{} \ill{} propriété (involutif ou pas) liste étant tel que
\[
  u(x) = x + \sigma(x) = 0 \quad \text{si } x \in E^{\circ}
\]
(exprimant que $\sigma | E^{\circ} = -\mathrm{id}_{E^{\circ}}$) donc $u$ se factorise en
\[
  \pi : F \longrightarrow E \qquad u = \pi \circ \varepsilon
\]
et \struck{$u(x)$} $\sigma$ s'exprime via $\pi$ comme
\[
  \underbrace{\sigma(x)}_{u(x)} = \pi(\varepsilon(x)) - x .
\]
Si on exprime maintenant que $\sigma$ induit $\mathrm{id}_F$ dans $E/E^{\circ}$, on trouve la condition
\[
  x \equiv \pi \varepsilon(x) - x \quad (E^{\circ}) \qquad \forall x \in E
\]
i.e.
\[
  \pi \underbrace{\varepsilon(x)}_{y} \equiv 2x \qquad \forall x \in E
\]
ce qui signifie aussi que la $\varepsilon$ des deux membres sont égaux, soit (posant $\varepsilon(x) = y$)
\[
  \boxed{\varepsilon \pi(y) = 2y}
\]
Ainsi, $\pi$ est « presque » une section, c'est une 2-section. \struck{Cela} Plus généralement, si $\lambda \in k$ et si on a une \uncertain{flèche}

\page{172}
\[
  \pi : F \longrightarrow E \quad \text{telle que} \quad \varepsilon \pi = \lambda\, \mathrm{id}_F ,
\]
cela signifie que l'extension ${}^{\lambda}E$ de $F$ par $E$, déduite de $E$ par l'homom.\ $\lambda\, \mathrm{id}_F$ de $F$, i.e.\ par la $\sigma$ \uncertain{induite} par $\lambda\, \mathrm{id}_E$ dans $E$, est munie d'un scindage.

On dit que $E$ est \emph{$\lambda$-scindée}.

Si p.~ex.\ $F = k$ (cas d'un torseur $A = \varepsilon^{-1}(1)$ sous $E^{\circ}$), la donnée de $\pi$ (\add{$: k \to E$}) revient à : la donnée d'un élément $t$ de $E$, et la condition envisagée c'est que
\[
  \varepsilon(t) = \lambda .
\]
Cela signifie aussi que le torseur
\[
  \varepsilon^{-1}(\lambda) = {}^{\lambda}A
\]
\struck{$\varepsilon^{-1}(A)$} déduit de $A$ par l'homomorphisme $\lambda\, \mathrm{id}_{A^{\circ} = E^{\circ}}$ sur le module des translations, est muni d'une section, « origine ». Dans l'explicitation
\[
  x \longmapsto \lambda x \qquad A \longrightarrow {}^{\lambda}A
\]
compatible avec $A^{\circ} \xrightarrow{\lambda\,\mathrm{id}} A^{\circ}$.
\note{dans $\varepsilon^{-1}(\lambda)$, l'argument est écrit par-dessus un $A$ ; « $\lambda\,\mathrm{id}_E$ » est lu sous réserve.}

\page{173}
peut être interprétée, par l'identification
\[
  {}^{\lambda}A \simeq A^{\circ}, \qquad x \longmapsto x - t
\]
comme une application de $A$ dans $A^{\circ}$, de façon précise en définissant
\[
  u : A \longrightarrow A^{\circ} \qquad u(x) = \lambda x - t
\]
application affine compatible avec
\[
  A^{\circ} \xrightarrow{\ \lambda\,\mathrm{id}_{A^{\circ}}\ } A^{\circ}
\]
et inversement une telle application \struck{\ill{}} \add{s'interprète comme une} \struck{applic.\ en} $A^{\circ}$-tels \uncertain{arrivées} ${}^{\lambda}A \simeq A^{\circ}$, \uncertain{donc} \add{i.e.\ choix d'une origine} \add{* dans ${}^{\lambda}A$}. Dans le cas où $\lambda = 2$, pour une involution de $F$ \ill{} sur \ill{} dans \ill{} dessus, on a donc
\[
  u(x) = 2x - t = -(\sigma(x) - x),
\]
c'est donc ce qui signifie précisément notre formule
\[
  \sigma(x) - x .
\]
\note{les ajouts interlinéaires de la seconde moitié de page, serrés au-dessus de la ligne, sont lus sous toute réserve.}

Or dans notre construction des $\ast A_i$, ce sont ces applications $u_i(x)$ : \struck{\ill{}} car ce qui revient au même, au signe près, les $\sigma_i(x_i) - x_i$ qui jouaient le rôle crucial. \struck{Il \ill{} s'agit pas} Cette application $u(x)$ étant affine, se prolonge linéairement en
\[
  E \xrightarrow{\ u\ } A^{\circ} \qquad u(x) = \lambda x - \underbrace{\varepsilon(x)\, t}_{= \pi \varepsilon(x)}
\]
\note{dans la dernière formule, $\lambda x$ est écrit par-dessus un autre symbole.}

\page{174}
Revenons au cas d'une ext.\ générale
\[
  0 \longrightarrow E^{\circ} \longrightarrow E \longrightarrow F \longrightarrow 0,
  \qquad F \simeq E/E^{\circ}
\]
et d'une $\lambda$-section
\[
  \pi : F \longrightarrow E ,
\]
associons lui
\begin{gather*}
  u = u^{\pi} : E \longrightarrow E^{\circ} \\
  u^{\pi}(x) = \lambda x - \pi \varepsilon(x)
\end{gather*}
\note{le premier $u$ est écrit par-dessus un $\pi$ biffé.}
Elle envoie bien $E$ dans $E^{\circ}$, car
\[
  \varepsilon u(x) = \lambda \underbrace{\varepsilon(x)}_{y} - \varepsilon\pi(\varepsilon(x))
  = \lambda y - \underbrace{\varepsilon\pi(y)}_{\lambda y} = 0
\]
d'autre part $u$ restreinte à $E^{\circ}$ est $\lambda\, \mathrm{id}_{E^{\circ}}$
\[
  u | E^{\circ} = \lambda\, \mathrm{id}_{E^{\circ}} \qquad \text{i.e. } u(x) = \lambda x \text{ si } x \in E^{\circ}
\]
C'est la version duale des $\lambda$-scindages. On récupère $\pi$ par $u$, par la formule
\[
  \pi(\varepsilon(x)) = \lambda x - u(x)
\]
en notant que le second membre se factorise bien par $F$, car il est nul sur $E^{\circ}$. La donnée de $\pi$ et de $u$ est équivalente. Mais c'est $u$ qui sera utilisé pour la définition des $\ast$.

\page{175}
On se fixe un $\lambda \in k$, et on travaille avec les extensions $\lambda$-scindées. \struck{\ill{} \ill{}} Si \ill{} de telles
\[
  0 \longrightarrow E_i^{\circ} \xrightarrow{\ i_{E_i}\ } E_i \longrightarrow F_i \longrightarrow 0
  \qquad u_i \circ i_{E_i} = \lambda\, \mathrm{id}_{E_i^{\circ}}
\]
\note{sous la flèche $E_i^{\circ} \to E_i$, une flèche courbe étiquetée $u_i$ revient de $E_i$ vers $E_i^{\circ}$ ; à droite de la formule, « $\lambda$ \uncertain{indépendant de $i$} », doublement souligné.}
et un module $B$, sur $k$. \struck{\ill{}}

On appelle $\lambda$-accouplement : à valeurs dans $B$, une couple
\[
  \textstyle\prod E_i \xrightarrow{\ f\ } B \qquad \prod E_i^{\circ} \xrightarrow{\ f^{\circ}\ } B
\]
d'applications multilinéaires, satisfaisant
\begin{align*}
  f(x_1, \ldots, x_i + \alpha_i, \ldots, x_n) &= f(x_1, \ldots, x_i, \ldots, x_n) \\
  &\quad + f^{\circ}(u_1(x_1), \ldots, \alpha_i, \ldots, u_n(x_n))
\end{align*}
Prenant les $x_i$ dans $E_i^{\circ}$, on trouve
\[
  f(\alpha_1, \ldots, \alpha_i, \ldots, \alpha_n) = \lambda^{n-1} f^{\circ}(\alpha_1, \ldots, \alpha_i, \ldots, \alpha_n)
\]
(on a pris $x_j = \alpha_j$ pour $j = 1, \ldots, \hat{\imath}, \ldots, n$, et $x_i = 0$). On \ill{} -- une solution du pb universel, notée
\[
  \mathop{\ast}_{i,\lambda} E_i \qquad (E_1 \mathbin{\ast_\lambda} E_2 \text{ etc.})
\]
l'indice $\lambda$ rappelant que c'est relatif à un choix de $\lambda$-trivialisations des extensions envisagées. L'expression $\mathop{\ast}_{i,\lambda} E_i$ se construit

\page{176}
comme quotient \struck{\ill{}} de
\[
  \textstyle\bigotimes_i E_i \oplus \bigotimes_i E_i^{\circ}
\]
et on a une suite exacte
\[
  \textstyle\bigotimes_i E_i^{\circ} \longrightarrow \mathop{\ast}_{i,\lambda} E_i \longrightarrow \bigotimes_i F_i \longrightarrow 0
\]
qui correspond au fait que les systèmes $(f, f^{\circ})$ pour lesquels $f^{\circ} = 0$, s'identifient aux $f$ qui se factorisent par les $F_i$. L'application $\bigotimes E_i^{\circ} \xrightarrow{\ i_\ast\ } \mathop{\ast}_{i,\lambda} E_i$ n'est sans doute pas toujours injective, mais je dis qu'elle l'est \uncertain{si} le choix des scindages des extensions $E_i^{\circ} \to E_i$ dépend \uncertain{d'un} \uncertain{nouveau} : \uncertain{générateur} des $\mathrm{ing}_i$.
\note{la fin de la phrase (« dépend … $\mathrm{ing}_i$ ») est d'une lecture très incertaine.}

En effet, écrivons
\[
  E_i \simeq E_i^{\circ} \oplus F_i
\]
la donnée des $u_i$, déjà connue sur $E_i^{\circ}$ (égale à $\lambda\,\mathrm{id}$) équivaut à la donnée des
\[
  t_i : F_i \longrightarrow E_i^{\circ}
\]
\note{le $t_i$ est écrit par-dessus un $b_i$.}

\page{177}
l'homomorphisme voulu de \struck{\ill{}} $E_1 \mathbin{\ast_\lambda} E_2$ dans $E_1^{\circ} \ast E_2^{\circ}$.

Pour \uncertain{préciser} : la linéarité \uncertain{étant} une condition \ill{} l'injectivité si les $F_i$ sont plats.

Il faut définir un $\lambda$-scindage \uncertain{canonique}, \uncertain{sur} $\mathop{\ast}_{i} E_i$, donc il faut définir
\[
  \mathop{\ast}_{i} E_i \longrightarrow \textstyle\bigotimes E_i^{\circ}
\]
qui sur $\bigotimes E_i^{\circ}$ \uncertain{donne} la valeur $\lambda\,\mathrm{id}$. Il faut donc trouver un $\lambda$-accouplement à valeurs dans $\bigotimes E_i^{\circ}$.

On \uncertain{a} \uncertain{tout} \uncertain{ce qu'il faut} :
\[
  f = \textstyle\bigotimes u_i : \bigotimes E_i \longrightarrow \bigotimes E_i^{\circ}
\]
\[
  f(x_1, \ldots, x_n) = u_1(x_1) \otimes \cdots \otimes u_n(x_n)
\]
\begin{align*}
  f(x_1 + \alpha_1, x_2, \ldots, x_n) &= (u_1(x_1) + \lambda\alpha_1) \otimes u_2(x_2) \otimes \cdots \otimes u_n(x_n) \\
  &= f(x_1, \ldots, x_n) + \underbrace{\lambda\alpha_1}_{\text{\struck{\ill{}}}} \otimes u_2(x_2) \otimes \cdots \otimes u_n(x_n)
\end{align*}
donne comme application
\[
  f^{\circ}(\alpha_1, \ldots, \alpha_n) = \lambda\, \alpha_1 \otimes \cdots \otimes \alpha_n
\]
\uncertain{du} \uncertain{moment} que les termes \ill{} \uncertain{précédents} \ill{}.
\note{la fin de la page est d'une lecture très incertaine.}
\marginal{mais il faudrait remplacer $\bigotimes E_i$ par un quotient -- on \uncertain{prend} bien un \ill{} dans les données $F_i \xrightarrow{\beta_i} E_i$ pas seulement injectives mais scindées, données d'un $u_i : E_i \to E_i^{\circ}$, avec $u_i \beta_i \equiv \lambda\,\mathrm{id}_{F_i}$ \ill{}}
\note{la note marginale est écrite en oblique dans le coin inférieur gauche du feuillet ; la seconde ligne de la formule développée porte, sous $\lambda\alpha_1$, un groupe raturé illisible ; la congruence finale de la note marginale est lue sous réserve.}

\page{178}
Ceci dit, pour définir un \add{\ill{}} $f$ (\struck{\ill{}} \uncertain{convenable}) pour un $f^{\circ}$ déjà donné, on peut se donner arbitrairement $g$ \add{(soit $g$)} sur $\bigotimes F_i$ (comme application multilin.)

\struck{On aura alors} Faisons le calcul pour deux facteurs.

On pose (notant $x_i \in F_i$, $\alpha_i \in E_i^{\circ}$)
\begin{align*}
  f(\underbrace{x_1 + \alpha_1}_{X_1}, \underbrace{x_2 + \alpha_2}_{X_2})
  &= g(x_1, x_2) + \text{\struck{\ill{}}}\, f^{\circ}(\alpha_1, b_2(x_2)) \\
  &\quad + f^{\circ}(b_1(x_1), \alpha_2) + \lambda \underbrace{f^{\circ}(\alpha_1, \alpha_2)}_{\substack{\shortparallel \\ f^{\circ}(u_1(\alpha_1), \alpha_2) \\ = f^{\circ}(\alpha_1, u_2(\alpha_2))}}
\end{align*}
\note{devant le premier $\lambda f^{\circ}(\alpha_1,\alpha_2)$ le $\lambda$ semble ajouté après coup ; sous $X_1$ le $X$ est repassé.}

Il suffit de vérifier que $f$ satisfait aux deux conditions. Or en remplaçant $\alpha_1$ par $\alpha_1 + \alpha'_1$, on obtient
\begin{align*}
  f(X_1 + \alpha'_1, X_2) &= f(X_1, X_2) + f^{\circ}(\alpha'_1, \underbrace{b_2(x_2)}) \\
  &\quad + \lambda f^{\circ}(\alpha'_1, \alpha_2)
\end{align*}
et on veut
\[
  = f(X_1, X_2) + f^{\circ}(\alpha'_1, \underbrace{u_2(X_2)}_{b_2(x_2) + \lambda\alpha_2})
  \qquad \text{OK.}
\]

Prenons pour $B$ le produit tensoriel $E_1^{\circ} \otimes E_2^{\circ}$, pour $f^{\circ}$ l'accouplement can., et $g = 0$, on trouve
\note{la page s'achève sur « on trouve » ; la suite est en tête de la page suivante.}

\page{179}
\[
  f(x_1, \ldots, x_n) = f^{\circ}(\alpha_1, u_2(x_2), \ldots, u_n(x_n)) .
\]
\note{la formule est reprise telle qu'écrite, avec $\alpha_1$ en premier argument ; un mot (« \uncertain{avec} ») précède $f^{\circ}$.}

On a $\ast_{i}$, \ill{} -- une opération interne $\ast$ dans la catégorie des extensions $\lambda$-scindées (à quotients $F_i$ plats, au besoin), et une particularité dans le \uncertain{tensoriel} \uncertain{gerbe} des \uncertain{torseurs} $\lambda$-scindés (cas où les $F_i$, sont $\simeq k$, $\bigotimes F_i \simeq k$ canon.!). Mais \uncertain{à} cela \ill{} -- \ill{} \uncertain{à} la commutativité, \uncertain{la} \uncertain{vérifier} \uncertain{lors} \uncertain{de} \uncertain{démontrer} (\uncertain{à vérifier} \uncertain{lemme} \uncertain{compatibilité} !), a-t-on une unité ?

Il faudrait pour ce \uncertain{faire}, \struck{\ill{}} pour une unité, \uncertain{prendre}
\[
  0 \longrightarrow E^{\circ} \longrightarrow E \longrightarrow F \longrightarrow 0,
\]
où $E^{\circ}$ et $F$ soient unités pour le $\otimes$ ordinaire, donc \struck{i\ill{}} $\simeq k$.
\[
  0 \longrightarrow k \longrightarrow E \longrightarrow k \longrightarrow 0
\]
\struck{\ill{}} Choisissons \uncertain{le} scindage, alors le $\lambda$-scindage \uncertain{sera} défini par un \struck{\ill{}} $t : k \xrightarrow{\ \beta\ } k$ i.e.\ $\beta \in k$. J'aimerais \uncertain{prendre} $\beta = 1$. Prenons plus généralement $\beta$ quelconque, je dis \uncertain{que}
\[
  E_\beta \mathbin{\ast_\lambda} E_2 \simeq \beta E_2 ,
\]
\uncertain{où} \ill{} $\beta E_2$ \uncertain{trivialisation} \uncertain{de} \uncertain{celle} de $\beta u_2$ \uncertain{déduite} de $E_2$.
\marginal{NB Si les ext.\ $E_i$ sont munies de scindages et donc de structures de $\lambda$-triv.\ définies par des $t_i : F_i \to E_i^{\circ}$, alors la $\lambda$-trivialisation de $\ast E_i$, scindée elle aussi, est définie par $t : \bigotimes F_i \to \bigotimes E_i^{\circ}$ \ill{}}
\note{la note marginale est écrite en oblique dans le coin supérieur gauche du feuillet ; sa fin, qui se perd dans le corps du texte, est illisible. Dans la ligne « donc $\simeq k$ », un mot biffé commençant par $i$ précède $\simeq k$.}

\page{180}
Ce faisant les $E \mapsto {}^{\beta}E$ \struck{\ill{}} sont des équivalences de cat.\ ssi $\beta$ est inversible dans $k$ -- et si $\beta = 1$, on retrouve ${}^{1}E = E$.

En termes intrinsèques, \struck{\ill{}} pour que \uncertain{les} $M$ $\lambda$-trivialisés, les conditions suivantes sont équivalentes
\begin{enumerate}
  \item $M \mathbin{\ast_\lambda} ?$ est une équivalence sur la cat.\ des $\lambda$-trivialisées
  \item $M^{\circ} \to M$ injectif, $\underset{L'}{M^{\circ}}$ et $\underset{L''}{M/M^{\circ}}$ inversibles, $M \xrightarrow{u} M^{\circ}$ est surjectif

  (en termes d'\struck{\ill{}} un scindage est quelconque que l'idéal \uncertain{donné} des \uncertain{valeurs} de $L' \otimes L^{-1} \xrightarrow{\beta} k$ et $\lambda k$ est $k$)
  \item $M$ inv.\ pour $\mathbin{\ast_\lambda}$
\end{enumerate}
\note{la liste s'arrête sur l'item 3 ; le bas de la page est blanc, et la question de l'unité pour $\ast_\lambda$ reste ouverte à la fin du lot. Les lectures de la parenthèse sous l'item 2 sont très incertaines ; dans $L' \otimes L^{-1}$, l'exposant de $L'$ est lu sous réserve.}

\end{document}

